\begin{helper}
Let \(W\) be a finite type, let \(H\) be a simple graph on \(W\), and let
\(s,k\) be natural numbers.  If \(H\) contains no \(s\)-clique and
\(H^{c}\) contains no \(k\)-clique, then
\(\mathrm{RamseyProperty}(s,k,|W|)\) is false.
\end{helper}
\begin{proof}
Let \(e:W\to \operatorname{Fin}(|W|)\) be the equivalence supplied by the
finite type structure, and define a graph \(G_{\operatorname{fin}}\) on
\(\operatorname{Fin}(|W|)\) by pulling \(H\) back along \(e^{-1}\).  Thus, for
vertices \(x,y\in \operatorname{Fin}(|W|)\), adjacency in
\(G_{\operatorname{fin}}\) is exactly adjacency of \(e^{-1}x\) and
\(e^{-1}y\) in \(H\).

Assume for contradiction that
\(\mathrm{RamseyProperty}(s,k,|W|)\) holds.  By
\(\noderef{RamseyProperty}\), applied to the graph \(G_{\operatorname{fin}}\),
there are two cases.

First suppose there is a finite set \(S\subseteq \operatorname{Fin}(|W|)\)
which is an \(s\)-clique in \(G_{\operatorname{fin}}\).  Map \(S\) into \(W\)
using the embedding \(e^{-1}\).  Since this map is injective, the image has
the same cardinality as \(S\), namely \(s\).  If two distinct vertices in the
image are chosen, they come from two distinct vertices of \(S\); because \(S\)
is a clique in \(G_{\operatorname{fin}}\), those two preimage vertices are
adjacent in \(G_{\operatorname{fin}}\), and by the definition of the pullback
graph their images under \(e^{-1}\) are adjacent in \(H\).  Hence the image is
an \(s\)-clique in \(H\), contradicting the hypothesis that \(H\) has no
\(s\)-clique.

Second suppose there is a finite set \(I\subseteq \operatorname{Fin}(|W|)\)
which is a \(k\)-clique in \(G_{\operatorname{fin}}^{c}\).  Again map \(I\)
into \(W\) using \(e^{-1}\).  Injectivity preserves its cardinality, so the
image has cardinality \(k\).  For two distinct vertices in the image, their
preimages in \(I\) are distinct.  Since \(I\) is a clique in the complement of
\(G_{\operatorname{fin}}\), the preimages are not adjacent in
\(G_{\operatorname{fin}}\); by the pullback definition, their images are not
adjacent in \(H\).  The same injectivity gives the vertices are distinct in
\(W\).  Therefore the image is a \(k\)-clique in \(H^{c}\), contradicting the
hypothesis that \(H^{c}\) has no \(k\)-clique.

Both alternatives in the Ramsey property lead to contradictions.  Therefore
\(\mathrm{RamseyProperty}(s,k,|W|)\) is false.
\end{proof}
